% GENERATED FILE. Edit canonical lessons, then run npm run generate:books.
% canonical-source-hash: fnv1a64:97ad54394c9d9dab
% canonical-lessons: SA-C141-pibati, SA-C141-vasati, SA-C141-carati, SA-C141-sprsati, SA-C141-tyajati

\chapter{Drinks, Dwells, Moves about, Touches, Leaves}
\label{ch:sa-drinks-dwells-moves-about-touches-leaves}
\hlchaptermodality{141}

\begin{chapteropening}
\textbf{By the end of this chapter:} \emph{I can say drinks, dwells, moves about, touches and leaves.}

Complete the last lesson of chapter 141: I can say drinks, dwells, moves about, touches and leaves.
\end{chapteropening}

\section[\texorpdfstring{\sk{पिबति} (pibati)}{पिबति (pibati)}]{\sk{पिबति} (pibati) — drinks}

\label{lesson:SA-C141-pibati}



\begin{quote}
Before the new one: say the Sanskrit for a bill (to pay), then the Sanskrit for a receipt.
\end{quote}

\subsection*{You'll want to know: \sk{पिबति}}
\textbf{\sk{पिबति}} — \emph{pibati} — \textquotedblleft{}drinks\textquotedblright{}.

\begin{cousinweb}
\emph{pibati} comes from the root √pā \textquotedblleft{}drink\textquotedblright{}, with its reduplicated present stem \emph{piba-} (MW). It is inherited from Indo-European: Latin \emph{bibo} \textquotedblleft{}I drink\textquotedblright{} and Greek \emph{pinō} are its cousins (EtymDB). It lives on as Hindi \emph{pīnā} \textquotedblleft{}to drink\textquotedblright{} (CDIAL 8209).
\end{cousinweb}

\begin{grammarlens}[title={What you've built}]
One more word for this chapter.
\end{grammarlens}

\subsection*{Your turn}
\begin{itemize}
\raggedright
  \item \emph{Say it:} \emph{pibati}
  \item \emph{Say it:} \emph{pibati}, once more
  \item \emph{Say it:} say \emph{pibati}
  \item \emph{Recall:} say the Sanskrit for a bill (to pay), then the Sanskrit for a receipt, then say \emph{pibati} again
\end{itemize}

\begin{tcolorbox}[breakable,colback=teal!4,colframe=teal!35!black,title={Before you move on}]
What does \sk{पिबति} mean? (\textquotedblleft{}Drinks\textquotedblright{}.) Say it once more.
\end{tcolorbox}
\section[\texorpdfstring{\sk{वसति} (vasati)}{वसति (vasati)}]{\sk{वसति} (vasati) — dwells, lives}

\label{lesson:SA-C141-vasati}



\begin{quote}
Before the new one: say the Sanskrit for a receipt, then the Sanskrit for drinks.
\end{quote}

\subsection*{You'll want to know: \sk{वसति}}
\textbf{\sk{वसति}} — \emph{vasati} — \textquotedblleft{}dwells, lives\textquotedblright{}.

\begin{cousinweb}
\emph{vasati} comes from the root √vas \textquotedblleft{}dwell, stay\textquotedblright{} (CDIAL 11435), inherited from Indo-European (EtymDB). MW sets it beside Gothic \emph{wisan} and English \emph{was} and \emph{were}: \textquotedblleft{}to stay\textquotedblright{} became \textquotedblleft{}to be\textquotedblright{}. It lives on as Hindi \emph{basnā} \textquotedblleft{}to dwell\textquotedblright{} (CDIAL 11435).
\end{cousinweb}

\begin{grammarlens}[title={What you've built}]
One more word for this chapter.
\end{grammarlens}

\subsection*{Your turn}
\begin{itemize}
\raggedright
  \item \emph{Say it:} \emph{vasati}
  \item \emph{Say it:} \emph{vasati}, once more
  \item \emph{Say it:} say \emph{vasati}
  \item \emph{Recall:} say the Sanskrit for a receipt, then the Sanskrit for drinks, then say \emph{vasati} again
\end{itemize}

\begin{tcolorbox}[breakable,colback=teal!4,colframe=teal!35!black,title={Before you move on}]
What does \sk{वसति} mean? (\textquotedblleft{}Dwells, lives\textquotedblright{}.) Say it once more.
\end{tcolorbox}
\section[\texorpdfstring{\sk{चरति} (carati)}{चरति (carati)}]{\sk{चरति} (carati) — moves about, grazes}

\label{lesson:SA-C141-carati}



\begin{quote}
Before the new one: say the Sanskrit for drinks, then the Sanskrit for dwells.
\end{quote}

\subsection*{You'll want to know: \sk{चरति}}
\textbf{\sk{चरति}} — \emph{carati} — \textquotedblleft{}moves about, grazes\textquotedblright{}.

\begin{cousinweb}
\emph{carati} comes from the root √car \textquotedblleft{}move, range about\textquotedblright{} (CDIAL 4686), inherited from Indo-European: Greek \emph{pelō} and Latin \emph{colō} are counted as its cousins (EtymDB). It lives on as Hindi \emph{carnā} \textquotedblleft{}to graze\textquotedblright{} (CDIAL 4686).
\end{cousinweb}

\begin{grammarlens}[title={What you've built}]
One more word for this chapter.
\end{grammarlens}

\subsection*{Your turn}
\begin{itemize}
\raggedright
  \item \emph{Say it:} \emph{carati}
  \item \emph{Say it:} \emph{carati}, once more
  \item \emph{Say it:} say \emph{carati}
  \item \emph{Recall:} say the Sanskrit for drinks, then the Sanskrit for dwells, then say \emph{carati} again
\end{itemize}

\begin{tcolorbox}[breakable,colback=teal!4,colframe=teal!35!black,title={Before you move on}]
What does \sk{चरति} mean? (\textquotedblleft{}Moves about, grazes\textquotedblright{}.) Say it once more.
\end{tcolorbox}
\section[\texorpdfstring{\sk{स्पृशति} (spṛśati)}{स्पृशति (spṛśati)}]{\sk{स्पृशति} (spṛśati) — touches}

\label{lesson:SA-C141-sprsati}



\begin{quote}
Before the new one: say the Sanskrit for dwells, then the Sanskrit for moves about.
\end{quote}

\subsection*{You'll want to know: \sk{स्पृशति}}
\textbf{\sk{स्पृशति}} — \emph{spṛśati} — \textquotedblleft{}touches\textquotedblright{}.

\begin{cousinweb}
\emph{spṛśati} comes from the root √spṛś \textquotedblleft{}touch, stroke\textquotedblright{} (CDIAL 13815), already in the Rigveda. \emph{sparśa} \textquotedblleft{}touch\textquotedblright{}, the word for the stop consonants in Sanskrit grammar, comes from the same root (MW).
\end{cousinweb}

\begin{grammarlens}[title={What you've built}]
One more word for this chapter.
\end{grammarlens}

\subsection*{Your turn}
\begin{itemize}
\raggedright
  \item \emph{Say it:} \emph{spṛśati}
  \item \emph{Say it:} \emph{spṛśati}, once more
  \item \emph{Say it:} say \emph{spṛśati}
  \item \emph{Recall:} say the Sanskrit for dwells, then the Sanskrit for moves about, then say \emph{spṛśati} again
\end{itemize}

\begin{tcolorbox}[breakable,colback=teal!4,colframe=teal!35!black,title={Before you move on}]
What does \sk{स्पृशति} mean? (\textquotedblleft{}Touches\textquotedblright{}.) Say it once more.
\end{tcolorbox}
\section[\texorpdfstring{\sk{त्यजति} (tyajati)}{त्यजति (tyajati)}]{\sk{त्यजति} (tyajati) — leaves, gives up}

\label{lesson:SA-C141-tyajati}



\begin{quote}
Before the new one: say the Sanskrit for moves about, then the Sanskrit for touches.
\end{quote}

\subsection*{You'll want to know: \sk{त्यजति}}
\textbf{\sk{त्यजति}} — \emph{tyajati} — \textquotedblleft{}leaves, gives up\textquotedblright{}.

\begin{cousinweb}
\emph{tyajati} comes from the root √tyaj \textquotedblleft{}leave, abandon, give up\textquotedblright{} (CDIAL 5984), inherited from Indo-European (EtymDB). \emph{tyāga} \textquotedblleft{}renunciation\textquotedblright{} comes from the same root.
\end{cousinweb}

\begin{grammarlens}[title={What you've built}]
One more word for this chapter.
\end{grammarlens}

\subsection*{Your turn}
\begin{itemize}
\raggedright
  \item \emph{Say it:} \emph{tyajati}
  \item \emph{Say it:} \emph{tyajati}, once more
  \item \emph{Say it:} say \emph{tyajati}
  \item \emph{Recall:} say the Sanskrit for moves about, then the Sanskrit for touches, then say \emph{tyajati} again
\end{itemize}

\begin{tcolorbox}[breakable,colback=teal!4,colframe=teal!35!black,title={Before you move on}]
What does \sk{त्यजति} mean? (\textquotedblleft{}Leaves, gives up\textquotedblright{}.) Say it once more.
\end{tcolorbox}
